\section{Key Stakeholders \& Blocs}

\subsection{Bloc Positions and Key Stakeholders}

The issue of international security in maritime chokepoints encompasses a rather diverse range of stakeholders whose strategic, economic, and security interests often converge yet conflict occasionally. While all states share a common interest in ensuring the uninterrupted flow of maritime commerce, their approaches differ significantly when analysed through the lens of sovereignty, freedom of navigation, military presence, and regional security. The following section outlines thoroughly the principal stakeholders and their respective positions concerning the five maritime chokepoints discussed in this Background Guide: the South China Sea, the Strait of Hormuz, the Bab el-Mandeb Strait, the Suez Canal and the Panama Canal.

\subsubsection{Permanent Members of the Security Council (P5)}

\emph{Countries Included}: USA, UK, France, Russia and China

The five permanent members of the Security Council play a decisive role in shaping international maritime security. Collectively, they possess significant naval capabilities, maintain global strategic interests, and influence the development and enforcement of international maritime law.

The United States, the United Kingdom and France generally advocate for a rules-based international order centred on freedom of navigation and unimpeded commercial shipping. They frequently support multinational naval operations aimed at protecting critical sea lanes.\footnote{Freund, E. (2017). Freedom of navigation in the South China Sea: A practical guide. Belfer Center for Science and International Affairs \url{https://www.belfercenter.org/publication/freedom-navigation-south-china-sea-practical-guide}} China, while equally dependent on secure maritime trade, places greater emphasis on sovereignty and territorial integrity, particularly in relation to the South China Sea\footnote{Wikipedia contributors. (n.d.). South China Sea Arbitration. In Wikipedia. Retrieved June 25, 2026, from \url{https://en.wikipedia.org/wiki/South_China_Sea_Arbitration}}. Russia similarly prioritizes state sovereignty and remains cautious of expanded multinational military operations near strategically sensitive maritime regions\footnote{Stronski, P. (2021, May). What is Russia doing in the Black Sea? Carnegie Endowment for International Peace. \url{https://carnegieendowment.org/posts/2021/05/what-is-russia-doing-in-the-black-sea}}.

Russia’s position is further shaped by its broader maritime strategy. Maintaining access to warm-water ports, preserving freedom of movement between the Black Sea and the Mediterranean, and limiting Western naval influence remain central to Russian security policy.\footnote{Stronski, P. (2021, May). What is Russia doing in the Black Sea? Carnegie Endowment for International Peace. \url{https://carnegieendowment.org/posts/2021/05/what-is-russia-doing-in-the-black-sea}} Since the outbreak of the conflict in Ukraine, issues surrounding maritime access, sanctions enforcement, and naval operations in adjacent seas have reinforced Moscow’s emphasis on sovereignty, strategic deterrence and opposition to what it perceives as external interference.\footnote{Chatham House. (2025, July). Understanding Russia's Black Sea strategy: Russia's global agenda through the Black Sea. \url{https://www.chathamhouse.org/2025/07/understanding-russias-black-sea-strategy}}

Given their political influence and military capabilities, the positions adopted by the P5 largely determine the Security Council’s ability to respond to maritime crises.

\subsubsection{Case Study: The South China Sea Disputes (2016 - Present)}

The South China Sea has become one of the world's most significant maritime security flashpoints, carrying approximately 1/3rd of global maritime trade annually\footnote{ChinaPower Project. (2017). How much trade transits the South China Sea? Center for Strategic and International Studies. \url{https://chinapower.csis.org/much-trade-transits-south-china-sea/}} while containing rich fisheries and potential hydrocarbon reserves. China claims sovereignty over most of the sea through its "Nine-Dash Line," which overlaps with the Exclusive Economic Zones (EEZs) claimed by several Southeast Asian states\footnote{Wikipedia contributors. (n.d.). South China Sea Arbitration. In Wikipedia. Retrieved June 25, 2026, from \url{https://en.wikipedia.org/wiki/South_China_Sea_Arbitration}}. Since the mid - 2010s, China has undertaken extensive land reclamation and militarization of artificial islands, constructing airstrips, ports, radar installations, and missile systems to strengthen its presence in the region\footnote{Newsweek. (2025, September 12). Satellite photos show Chinese military outposts in the South China Sea. \url{https://www.newsweek.com/satellite-photos-show-chinese-military-outpost-south-china-sea-2128785}}.

In 2016, the Permanent Court of Arbitration ruled that China's expansive maritime claims had no legal basis under UNCLOS\footnote{Permanent Court of Arbitration. (2016, July 12). The South China Sea Arbitration (The Republic of the Philippines v. The People's Republic of China), PCA Case No. 2013-19. \url{https://pca-cpa.org/en/cases/7/}}, a decision China blatantly rejected. In response, the United States has continued conducting Freedom of Navigation Operations (FONOPs) to challenge excessive maritime claims and uphold what it considers the rules-based international order\footnote{Poling, G. (2025, July 22). The 2016 South China Sea arbitration and the limits of international law. The Diplomat.}. The United Kingdom and France have similarly supported freedom of navigation through naval deployments and diplomatic statements, while Russia has generally supported China's emphasis on sovereignty and opposed external military involvement in regional disputes\footnote{Asia Maritime Transparency Initiative. (2016, July 11). Judgment day: The South China Sea tribunal issues its ruling. Center for Strategic and International Studies. \url{https://amti.csis.org/arbitration-ruling-analysis/}}.

This case vividly illustrates the fundamental divide amongst the P5 regarding the balance between territorial sovereignty, international law and the protection of global maritime commons.

\minorheading{Gulf and Red Sea Stakeholders}

\emph{Countries Included}: The Kingdom of Bahrain, the Federal Republic of Somalia, the Islamic Republic of Iran, the Kingdom of Saudi Arabia, the Arab Republic of Egypt, and the State of Israel

The Gulf and Red Sea region constitutes one of the world's most strategically sensitive maritime corridors, encompassing the Strait of Hormuz, the Bab el-Mandeb Strait and the Suez Canal. States within this region rely heavily on secure maritime routes for energy exports, international trade as well as regional stability.

Iran views the Strait of Hormuz as a vital component of its national security and frequently opposes an extensive foreign military presence in the Gulf\footnote{Iran International. (2026, July 4). Iran warns against foreign military presence in Strait of Hormuz after UK-France statement. \url{https://www.iranintl.com/en/202607042713}}. Conversely, Bahrain and Saudi Arabia support international cooperation aimed at protecting maritime commerce and ensuring uninterrupted energy exports. Egypt, as the primary administrator of the Suez Canal, prioritizes the security and operational continuity of one of the world's busiest waterways, whereas Israel remains particularly concerned with maintaining access through the Red Sea in light of recent attacks on its banks for commercial shipping linked to it\footnote{Wikipedia contributors. (n.d.). Houthi attacks on commercial vessels. In Wikipedia. Retrieved June 25, 2026, from \url{https://en.wikipedia.org/wiki/Houthi_attacks_on_commercial_vessels}} . Somalia, although not directly controlling a chokepoint, remains a critical stakeholder due to the continued threat of piracy and maritime insecurity within the Gulf of Aden and surrounding waters\footnote{United Nations Press. (2014, October 24). Security Council renews action to fight piracy off Somali coast, calling for deployment of vessels, arms, military aircraft (SC/11643). \url{https://press.un.org/en/2014/sc11643.doc.htm}}.

Despite differing geopolitical interests, these actors share a common interest in preventing prolonged disruptions to international shipping and regional economic stability.

\emph{Case Study: Houthi Attacks in the Bab el-Mandeb Strait (2023 - Present)}

The Bab el-Mandeb Strait serves as the southern gateway to the Red Sea and the Suez Canal, making it one of the world's most strategically important maritime chokepoints. Since late 2023, the security environment has deteriorated considerably as Yemen's Houthi movement launched repeated missile, drone and unmanned surface vessel attacks against commercial shipping and naval vessels transiting the region. These attacks were framed by the Houthis as support for Palestinians during the Gaza conflict but quickly evolved into a broader challenge to international maritime security\footnote{Atlas Institute for International Affairs. (2025, March 27). The Red Sea shipping crisis (2024–2025): Houthi attacks and global trade disruption. \url{https://atlasinstitute.org/the-red-sea-shipping-crisis-2024-2025-houthi-attacks-and-global-trade-disruption/}}.

The disruption prompted many global shipping companies to suspend Red Sea transit and reroute vessels around the Cape of Good Hope, adding thousands of nautical miles to voyages between Europe and Asia\footnote{S\&P Global Commodity Insights. (2023, December 21). FACTBOX: Seaborne trade reroutes away from Red Sea over Houthi attacks. \url{https://www.spglobal.com/commodityinsights/en/market-insights/latest-news/shipping/122123-factbox-seaborn} e-trade-reroutes-away-from-red-sea-over-houthi-attacks}. This significantly increased transportation costs, insurance premiums, fuel consumption and delivery times, affecting global supply chains and international trade\footnote{Atlas Institute for International Affairs. (2025, March 27). The Red Sea shipping crisis (2024–2025): Houthi attacks and global trade disruption.}. In response, the United States, the United Kingdom, and several partner states launched multinational maritime operations to protect commercial shipping, while regional actors adopted differing positions on military intervention and escalation\footnote{Atlas Institute for International Affairs. (2025, March 27). The Red Sea shipping crisis (2024–2025): Houthi attacks and global trade disruption.}.

For Gulf and Red Sea stakeholders, this crisis demonstrates the close relationship between regional conflicts and global maritime commerce. While all states share an interest in maintaining secure shipping lanes, they still remain divided over the role of foreign naval forces, regional security arrangements and the rooted, underlying political causes of maritime instability.

\minorheading{Indo-Pacific Maritime Stakeholders}

\emph{Countries Included}: China, Pakistan, India and Singapore

The Indo-Pacific region has become increasingly central to global maritime security, particularly due to the strategic significance of the South China Sea. This region facilitates a substantial proportion of global trade while simultaneously serving as an arena of geopolitical competition.

China regards the South China Sea as integral to its sovereignty and national security, whereas India seeks to preserve stability throughout the Indian Ocean and safeguard vital sea lines of communication connecting Asia with global markets. Pakistan similarly depends upon secure maritime trade routes through the Arabian Sea and Indian Ocean. Singapore, situated adjacent to one of the world's busiest shipping corridors, consistently advocates for open and secure navigation and supports multilateral approaches to maritime governance.

Although these states differ in their strategic priorities, they share an interest in preventing disruptions that could undermine regional stability and international commerce.

\subsubsection{Case Study: The 2016 Permanent Court of Arbitration Ruling}

In July 2016, the Permanent Court of Arbitration in The Hague issued a landmark ruling in the case brought by the Philippines against China under the United Nations Convention on the Law of the Sea (UNCLOS)\footnote{Permanent Court of Arbitration. (2016, July 12). The South China Sea Arbitration (The Republic of the Philippines v. The People's Republic of China), PCA Case No. 2013-19. \url{https://pca-cpa.org/en/cases/7/}}. The tribunal concluded that China's "Nine-Dash Line" claims lacked legal basis and affirmed that several disputed maritime features were not entitled to extensive maritime zones\footnote{Wikipedia contributors. (n.d.). South China Sea Arbitration. In Wikipedia. Retrieved June 25, 2026, from \url{https://en.wikipedia.org/wiki/South_China_Sea_Arbitration}}. Although legally significant, China rejected the ruling, arguing that the tribunal lacked jurisdiction and maintaining it’s historical claims over much of the South China Sea\footnote{Poling, G. (2025, July 22). The 2016 South China Sea arbitration and the limits of international law. The Diplomat. \url{https://thediplomat.com/2025/07/the-2016-south-china-sea-arbitration-and-the-limits-of-international-law/}}.

The ruling has since become a defining reference point for maritime governance in the Indo-Pacific. States such as Singapore and India have consistently emphasized the importance of international law, peaceful dispute resolution as well as the freedom of navigation, while Pakistan has generally aligned more closely with China's regional position. Meanwhile, China has continued strengthening its administrative and military presence across disputed features, reinforcing its sovereignty claims despite international criticism\footnote{Poling, G. (2025, July 22). The 2016 South China Sea arbitration and the limits of international law. The Diplomat. \url{https://thediplomat.com/2025/07/the-2016-south-china-sea-arbitration-and-the-limits-of-international-law/}}.

This case quietly reflects the competing interpretations of international maritime law and sovereignty that characterize the Indo-Pacific, as well as , puts a spotlight on the challenge of enforcing international legal decisions when major powers reject their legitimacy, making the South China Sea , yet again, one of the defining maritime security issues of the twenty-first century.

\minorheading{Canal States and relevant Global Trade}

\emph{Countries Included}: Arab Republic of Egypt, Republic of Panama

The Suez Canal and the Panama Canal are among the world's most critical artificial waterways, linking major oceans and significantly reducing transit times for international shipping. Henceforth, Egypt and Panama occupy unique positions as primary custodians of maritime infrastructure with such global economic and strategic importance.

Both states prioritize the uninterrupted operation, physical security and long-term resilience of their respective canals while maintaining sovereign authority over canal management. Egypt's concerns are closely tied to regional instability affecting the Red Sea and Suez route, whereas Panama faces growing challenges from climate change, water scarcity, and major-power interest in canal operations\footnote{Woodwell Climate Research Center. (2024). Drought, climate, and the Panama Canal. \url{https://www.woodwellclimate.org/drought-panama-canal-7-graphics/}}.

Recent disruptions, including the grounding of the \emph{Ever Given} in the Suez Canal\footnote{SAFETY4SEA. (2026, January 22). Ever Given: The grounding that changed the world's view of shipping. \url{https://safety4sea.com/cm-ever-given-the-grounding-that-changed-the-worlds-view-of-shipping/}} and drought-related transit restrictions in the Panama Canal\footnote{Woodwell Climate Research Center. (2024). Drought, climate, and the Panama Canal. \url{https://www.woodwellclimate.org/drought-panama-canal-7-graphics/}}, have exposed the vulnerability of global supply chains to both security and environmental risks. As a result, Egypt and Panama generally support international cooperation, infrastructure investment and resilience planning, aimed at keeping commercial shipping secure, efficient and politically neutral

\subsubsection{Case Study: The Ever Given Grounding in the Suez Canal (2021)}

In March 2021, the ultra-large container vessel \emph{Ever Given} became lodged diagonally across the Suez Canal after being caught in strong winds during transit. The blockage lasted six days and prevented hundreds of vessels from passing through one of the world's busiest maritime trade routes\footnote{SAFETY4SEA. (2026, January 22). Ever Given: The grounding that changed the world's view of shipping. \url{https://safety4sea.com/cm-ever-given-the-grounding-that-changed-the-worlds-view-of-shipping/}}. As approximately 12\% of global trade normally transits the canal, the incident quickly disrupted international supply chains, delayed energy shipments, and generated significant financial losses for shipping companies and global markets\footnote{New Zealand Ministry of Foreign Affairs and Trade. (2021, April 18). The importance of the Suez Canal to global trade. \url{https://www.mfat.govt.nz/en/trade/mfat-market-reports/the-importance-of-the-suez-canal-to-global-trade-18-apri} l-2021}.

Although the blockage resulted from an operational accident rather than conflict, it demonstrated the vulnerability of critical maritime infrastructure to even short-term disruptions. The incident prompted renewed international discussions regarding infrastructure resilience, emergency response capabilities, and the need to diversify global supply chains. Egypt subsequently invested in expanding sections of the canal and modernizing navigation systems to reduce the likelihood of similar incidents.

For Egypt, the incident reaffirmed the importance of safeguarding the Suez Canal against operational and security disruptions, whereas for Panama, it highlighted the broader vulnerability of global trade to disruptions in critical maritime infrastructure and reinforced the importance of maintaining a resilient, climate-adaptive and politically neutral canal capable of supporting uninterrupted international commerce.

\minorheading{Global Shipping Nations}

\emph{Countries Included}: The Hellenic Republic of Greece, the Kingdom of Denmark, the Republic of Liberia and the Republic of Singapore

Several states, despite lacking direct control over major maritime chokepoints, possess substantial interests in preserving secure international shipping. Greece, Denmark and Liberia collectively represent a significant proportion of the global merchant fleet through ship ownership and vessel registration\footnote{SAFETY4SEA. (2025, October 2). UNCTAD: The top registries and ship owning nations of the world. \url{https://safety4sea.com/unctad-the-top-registries-and-ship-owning-nations-of-the-world/}}.

Their economic prosperity is closely tied to the uninterrupted movement of commercial shipping. Consequently, these states consistently support freedom of navigation, secure maritime trade routes, effective counter-piracy measures, and international cooperation aimed at reducing risks to commercial vessels.

\subsubsection{Case Study: Red Sea Shipping Rerouting (2023 - Present)}

The escalation of attacks against commercial vessels in the Red Sea has significantly affected global maritime commerce. As security risks increased, many international shipping companies, including several of the world's largest container carriers, temporarily suspended transit through the Bab el-Mandeb Strait and the Suez Canal, opting instead to re-route vessels around the Cape of Good Hope\footnote{SAFETY4SEA. (2026, January 22). Ever Given: The grounding that changed the world's view of shipping. \url{https://safety4sea.com/cm-ever-given-the-grounding-that-changed-the-worlds-view-of-shipping/}}. While improving vessel safety, this alternative route substantially increased voyage durations, operational costs, fuel consumption and greenhouse gas emissions\footnote{S\&P Global Commodity Insights. (2023, December 21). FACTBOX: Seaborne trade reroutes away from Red Sea over Houthi attacks. \url{https://www.spglobal.com/commodityinsights/en/market-insights/latest-news/shipping/122123-factbox-seaborn} e-trade-reroutes-away-from-red-sea-over-houthi-attacks}.

These disruptions had immediate consequences for states whose economies rely heavily on international shipping, merchant fleets and maritime logistics. Greece, Denmark and Liberia, experienced increased insurance costs and operational uncertainty\footnote{S\&P Global Commodity Insights. (2023, December 21). FACTBOX: Seaborne trade reroutes away from Red Sea over Houthi attacks. \url{https://www.spglobal.com/commodityinsights/en/market-insights/latest-news/shipping/122123-factbox-seaborn} e-trade-reroutes-away-from-red-sea-over-houthi-attacks}. Even Singapore faced challenges due to disruptions affecting wider Asian trade networks.

\minorheading{Emerging Maritime Stakeholders}

\emph{Countries Included}: Republic of Latvia, Democratic Republic of the Congo and Republic of Türkiye

States such as Colombia, the Democratic Republic of the Congo and Latvia possess comparatively limited direct influence over the maritime chokepoints examined within this Background Guide. Nevertheless, they remain important participants within the Security Council due to the global consequences of maritime disruptions.

These states generally advocate for multilateral cooperation, adherence to international law, capacity building and the peaceful resolution of disputes, whilst their perspectives frequently emphasize strengthening international institutions, enhancing maritime governance and improving resilience against emerging security threats.

\subsubsection{Case Study: International Counter-Piracy Operations off the Coast of Somalia}

During the late 2000s and early 2010s, piracy along the Somali coast became one of the most serious non-state threats to maritime security. Armed groups targeted commercial vessels transiting the Gulf of Aden and the western Indian Ocean, endangering a major trade corridor between Europe, Asia and the Middle East. The international response combined coordinated naval patrols, information - sharing mechanisms, capacity-building efforts and Security Council resolutions authorizing counter - piracy measures\footnote{S\&P Global Commodity Insights. (2023, December 21). FACTBOX: Seaborne trade reroutes away from Red Sea over Houthi attacks. \url{https://www.spglobal.com/commodityinsights/en/market-insights/latest-news/shipping/122123-factbox-seaborn} e-trade-reroutes-away-from-red-sea-over-houthi-attacks}.

These efforts significantly reduced successful pirate attacks and demonstrated the value of multilateral cooperation in protecting civilian shipping. Unlike many territorial disputes, counter-piracy operations received broad international support because they addressed a shared security threat while strengthening regional maritime governance\footnote{ISS Africa. (2024). Is Somali piracy finally under control? Institute for Security Studies. \url{https://issafrica.org/iss-today/is-somali-piracy-finally-under-control}}.

Nevertheless, Türkiye offers a related yet distinct perspective. Through its control of the Bosporus and Dardanelles under the Montreux Convention of 1936,\footnote{Wikipedia contributors. (n.d.). Montreux Convention Regarding the Regime of the Straits. In Wikipedia. Retrieved June 26, 2026, from \url{https://en.wikipedia.org/wiki/Montreux_Convention_Regarding_the_Regime_of_the_Straits}} Türkiye has experience balancing coastal state sovereignty with international navigation.\footnote{Encyclopaedia Britannica. (2023). Montreux Convention. In Encyclopaedia Britannica. \url{https://www.britannica.com/event/Montreux-Convention}} Although these straits are not the focus of this agenda, Türkiye supports solutions respecting sovereign rights, promoting regional stability and avoiding unnecessary military escalation.
